%% ma: L12-B06-0090
%% lop: 12
%% chuong: 2
%% bai: 6
%% dang: 12.6.8
%% loai: DS
%% muc: [TH, TH, NB, TH]
%% dap_an: "ĐSSĐ"
%% nguon: "Ảnh giáo viên gửi 05/11/2026 (M11): Kiểm tra giữa HK1 Toán 12 (2026-2027), THPT Hồ Thị Bi (TP.HCM), mã đề 1133, Phần II Câu 2"
%% kien_thuc: "Bài 6 (tích vô hướng, góc giữa hai vectơ, độ dài vectơ, quy tắc hình bình hành); Lớp 11 (đường thẳng vuông góc với mặt phẳng)"
%% trang_thai: chinh-thuc
%% ghi_chu: "Đề gốc không kèm đáp án; đã giải độc lập. Mức từng ý: a) TH, b) TH, c) NB, d) TH. Hình dựng lại từ dữ kiện đề (SA vuông góc đáy), SA vẽ theo tỉ lệ a√2."
\begin{ex}
Cho hình chóp $S.ABCD$ có đáy $ABCD$ là hình vuông cạnh $a$, $SA\perp(ABCD)$; $SA=a\sqrt2$ (tham khảo hình vẽ).
\begin{center}
\begin{tikzpicture}[x={(1.8cm,0cm)},y={(.9cm,.6cm)},z={(0cm,1.5cm)}]
  \coordinate (A) at (0,1,0); \coordinate (B) at (1,1,0); \coordinate (C) at (1,0,0); \coordinate (D) at (0,0,0);
  \coordinate (S) at (0,1,1.414);
  \draw[khuat] (A)--(B) (A)--(D) (A)--(S);
  \draw (B)--(C)--(D) (S)--(B) (S)--(C) (S)--(D);
  \path (S) node[above]{$S$} (A) node[above left,font=\footnotesize]{$A$} (B) node[right]{$B$} (C) node[below]{$C$} (D) node[below left]{$D$};
\end{tikzpicture}
\end{center}
\choiceTF
{\True $\overrightarrow{SB}\cdot\overrightarrow{DC}=a^2$.}
{Số đo góc giữa hai vectơ $\overrightarrow{AB}$ và $\overrightarrow{CD}$ bằng $0^\circ$.}
{$\left|\overrightarrow{SB}\right|=a\sqrt5$.}
{\True $\overrightarrow{SA}+\overrightarrow{SC}=\overrightarrow{SB}+\overrightarrow{SD}$.}
\loigiai{a) $\overrightarrow{DC}=\overrightarrow{AB}$ nên $\overrightarrow{SB}\cdot\overrightarrow{DC}=(\overrightarrow{SA}+\overrightarrow{AB})\cdot\overrightarrow{AB}=\overrightarrow{SA}\cdot\overrightarrow{AB}+AB^2=0+a^2=a^2$ (vì $SA\perp AB$). Đúng.

b) $\overrightarrow{CD}=-\overrightarrow{AB}$ nên góc giữa $\overrightarrow{AB}$ và $\overrightarrow{CD}$ bằng $180^\circ$. Sai.

c) $SB^2=SA^2+AB^2=2a^2+a^2=3a^2$ nên $|\overrightarrow{SB}|=a\sqrt3\ne a\sqrt5$. Sai.

d) Gọi $O$ là tâm hình vuông $ABCD$: $\overrightarrow{SA}+\overrightarrow{SC}=2\overrightarrow{SO}=\overrightarrow{SB}+\overrightarrow{SD}$. Đúng.}
\end{ex}
